We identify the OPERATOR at precisely BFJ's published setting: a smooth CLOSED point of an affine variety over an algebraically closed characteristic-zero field, dimension \(n\ge2\), initially with gauge \(\mathfrak m\). G331 grants the full cone specialization on all semivaluations, with bottom adjoined; bottom is not bounded. No smooth multiplier statement is applied to singular or arbitrary-characteristic rings.

The complete BFJ Section 4 was read, not only the sequence statement 4.8. Its Cartier pairing is symmetric multilinear and has the ordinary exceptional intersection on a determination. Equations (4.3), Propositions 4.6--4.7 give ideal moments using the generic parameter COLENGTH convention. This equals our polynomial mixed multiplicity in the regular local ring, as follows.

A superficial parameter tuple exists by consecutive 17.2.2 choices and 17.3.1's joint-reduction construction. Joint reductions form a nonempty Zariski open set of coefficient tuples: for a fixed sufficiently large diagonal multi-Rees degree, the joint-reduction equation is surjectivity of a finite-dimensional multiplication matrix whose entries are linear in these coefficients, an open rank condition. Nakayama lifts residue surjectivity to the ideal equation. Conversely any joint reduction satisfies this equation in some such degree. Complete 17.4.9 says every joint reduction's parameter ideal has multiplicity equal to the polynomial mixed coefficient, at arbitrary Noetherian hypotheses. We use this ordinary published application here after reading its full statement/proof; it is not needed for the general monotonicity proof. In a regular local ring the parameter tuple is a regular sequence, so its Samuel multiplicity is exactly its quotient length. Thus generic parameter colength equals our mixed coefficient. Equivalently, for a superficial tuple repeated 17.4.6 has no annihilator correction by the CM property. These arguments are confined to the smooth range.

Hence BFJ 4.7 gives exactly n2's moments for rational primary profiles, with the same \(p=n-1\) arity and sign. Scaling, unit/zero conventions and lattice density give equality as Radon measures. For arbitrary \(C\) tuples, BFJ 4.4's smooth addition/order laws give the same negative-constant shift estimate (U). The mass is bounded by \(M^pe(\mathfrak m)\); uniform rational approximation consequently identifies its moments and measure with ours.

For bounded \(D\) tuples choose any actual decreasing \(C\)-representations. At each divisorial valuation their associated nef Weil coefficients decrease to the coefficients of the G331-identified BFJ limit, so the product net converges in the coefficient topology. BFJ 4.4's FULL proof gives upper semicontinuity in this topology: the strict coefficient inequality below a fixed nef Cartier divisor is an open finite-determination condition; Cartier majorants nearly attaining the infimum can be scaled slightly toward zero to make the inequalities strict. This provides open neighborhoods with intersection below any prescribed larger value. Together with intersection monotonicity, it proves continuity along ANY decreasing coefficient net: the values stay above the limit by order, while upper semicontinuity makes their limsup at most the limit. It needs no cofinal sequence. Apply this with the fixed ideal Cartier test \(Z(J)\). Positive ideal moments change the sign, and they converge exactly to our scalar limits in n2. Mass bounds and lattice density identify the entire measures on all semivaluations. The source's 4.8 sequence assertion alone was not promoted to a net theorem.

This is exact published smooth operator specialization, with the actual proof-topology argument. Remark 4.5's essential multiplier approximation for additivity remains a disclosed smooth dependency, as do resolution, Rees/Ramanujam intersections and usual CM multiplicity foundations. It supplies no general singular theorem. For other primary gauges use our proved (GAUGE); we do not attribute an unstated gauge theorem to BFJ. Smooth DVR and empty dimension zero follow directly from n6.
